% Standalone supplementary document. Compile after main.tex, or \input the two section files
% into main.tex before \end{document} to produce a single combined PDF.
\documentclass[twocolumn,10pt]{article}

\usepackage{microtype}
\usepackage{graphicx}
\usepackage{booktabs}
\usepackage{amsmath,amssymb}
\usepackage{xcolor}
\usepackage[hidelinks]{hyperref}
\usepackage{url}
\usepackage{caption}
\usepackage{natbib}
\usepackage[margin=0.75in]{geometry}

\IfFileExists{icml2025.sty}{\usepackage[accepted]{icml2025}}{}

\captionsetup{font=small,labelfont=bf,skip=4pt}
\setlength{\tabcolsep}{4pt}
\renewcommand{\arraystretch}{1.06}

\newcommand{\Hstar}{H^{*}}
\newcommand{\Hhat}{\hat{H}}
\newcommand{\gphi}{g_{\phi}}
\providecommand{\tabledir}{../tables}
\graphicspath{{../}}

% Cross-references to the main text resolve only in the combined build; standalone builds
% will show them as bold ?? markers, which is expected.
\title{Supplementary Material\\\large Learning to Monitor the Reliability Horizon of
Imagined Futures}
\author{Abhishek Rajdhar Dubey \\
  Independent Researcher \\
  \texttt{abhishekrajdhardubey@gmail.com}}
\date{}

\begin{document}
\maketitle

This supplement documents the environments and data collection, the world models and monitor,
the statistical procedure, the censoring analysis, additional results including the
size-matched control and the ensemble-size ablation, the two configurations rejected as invalid
instruments, the planning experiment, and the executable verification of the information
boundary.

Pre-registration documents, the deviation ledger and per-seed statistics are available
with the code at \url{https://github.com/abhishekrajdhar/reliability_in_worldmodel}.
Every number reported in the main paper and in this supplement is emitted programmatically from
stored result artifacts by \texttt{scripts/make\_paper\_tables.py} and
\texttt{scripts/make\_paper\_figures.py}; neither script recomputes an experiment.

\section{Experimental Details}
\label{app:details}

\subsection{Environments and data collection}

All configurations are MiniGrid environments \citep{chevalier2023minigrid} with
\texttt{agent\_view\_size} $=5$, giving a $5\times5\times3$ egocentric symbolic observation, and
\texttt{max\_steps} $=400$. Because the view is egocentric, grid size never enters the
observation and one world model applies to every configuration by construction.

\begin{center}\small
\begin{tabular}{llll}
\toprule
& base environment & perturbation & role \\
\midrule
A  & \texttt{DoorKey-16x16}        & ---                       & training distribution \\
B  & \texttt{FourRooms}            & ---                       & unseen layout topology \\
C  & \texttt{DoorKey-16x16}        & action slip $0.15$        & stochastic dynamics \\
D  & \texttt{MultiRoom-N6}         & ---                       & long corridors \\
E  & \texttt{SimpleCrossingS11N5}  & ---                       & unseen maze family \\
F  & \texttt{DoorKey-16x16}        & obs.\ noise $0.15$        & rejected instrument \\
F2 & \texttt{DoorKey-16x16}        & obs.\ noise $0.70$        & rejected instrument \\
F3 & \texttt{DoorKey-16x16}        & action dropout $0.20$     & accepted instrument \\
\bottomrule
\end{tabular}
\end{center}

\emph{Action slip} replaces the requested action by a uniformly random one. \emph{Observation
noise} corrupts a fraction of object-channel cells in the egocentric view, leaving dynamics
untouched. \emph{Action dropout} replaces the requested action by MiniGrid's \texttt{done}
action, a verified no-op that does not terminate the episode, so the agent's intent silently
fails to execute.

\paragraph{Collection policy.} Trajectories are collected by a forward-biased explorer
($p(\text{forward}) = 0.6$) that additionally pursues scripted key/door subgoals with
probability $0.35$, for $120$ steps per episode. An earlier BFS-based explorer was discarded
during pre-registration because it produced $P(\text{moved} \mid \text{forward}) = 0.985$:
shortest paths never walk into walls, which removes the epistemic uncertainty the experiment
measures. Under the forward-biased policy that quantity is $0.48$--$0.67$ across
configurations.

\paragraph{Splits.} Layout seeds are disjoint by construction: world-model training
$[0,3000)$, monitor training $[3000,4200)$, validation $[4200,4800)$, test $[4800,5400)$, and
shifted configurations $[100000,100600)$. The monitor trains on episodes the world model never
saw, so $\Hstar$ is not inflated by memorisation. For the transfer experiments each shifted
configuration is further partitioned by episode id ($\bmod\,5 \in \{0,1,2\}$ train, $3$
validation, $4$ test), with the test portion evaluation-only in every regime.

\subsection{World model}

Encoder: per-cell embeddings of object ($11 \to 16$), colour ($6 \to 6$) and state
($3 \to 6$), flattened over the $5\times5$ view and concatenated with a one-hot direction and a
carrying flag, then a two-layer MLP to $384$ units. Recurrent core: a GRU cell
\citep{cho2014gru} with deterministic state $512$ and Gaussian stochastic state $48$, with
prior $p(z_t \mid h_t)$ and posterior $q(z_t \mid h_t, e_t)$. Decoder heads on $[h_t, z_t]$
predict $\mathrm{dir}$ (4-way), $\mathrm{moved}$, $\mathrm{carrying}$ and
$\mathrm{door\_open}$. Approximately $2.85$M parameters.

Training deliberately mixes filtering and open-loop imagination: for each sequence a split
point $c \sim \mathcal{U}\{5,\dots,20\}$ is drawn, the posterior is used for $t < c$ and the
prior alone for $t \ge c$, so the model is optimised for the open-loop rollout that is later
evaluated. Loss is a weighted cross-entropy over the heads (weights $1.0$ for $\mathrm{dir}$,
$2.0$ for $\mathrm{moved}$, $0.3$ for the remaining two) plus a free-nats KL term.
Optimisation: AdamW, learning rate $3\times10^{-4}$, weight decay $10^{-5}$, batch $64$,
sequence length $50$, $6000$ steps, gradient-norm clipping at $100$, mixed precision on GPU.
Ensembles use $K=5$ members differing in seed and data order.

\paragraph{Transformer variant.} The Transformer world model replaces the recurrence with a
causal self-attention stack \citep{vaswani2017attention} ($d_{\text{model}} = 256$, $2$ layers, $4$ heads, GELU,
pre-norm). Observations are replaced by a learned mask token wherever unavailable, so training
and imagination are single parallel passes and there is no autoregressive latent feedback. It
exposes the same interface and head set. One definitional difference is unavoidable: its
one-step context error uses a fully masked pass, because a per-token ``everything strictly
before $t$'' mask is not expressible in one forward pass. This withholds \emph{more}
information than the GRU version, so it cannot flatter the Transformer.

\subsection{Monitor}

$\gphi$ is an MLP with two hidden layers of $256$ units and SiLU activations, producing logits
over $\Hstar \in \{0,\dots,30\}$, trained with cross-entropy for $60$ epochs (AdamW, learning
rate $10^{-3}$, weight decay $10^{-4}$, batch $512$). Inputs are standardised using statistics
from the training split only and clipped to $\pm 10$ s.d. The point estimate is the mean of the
predictive distribution, and the same convention is applied to every baseline including
isotonic regression, so comparisons are like-for-like. Checkpoint selection uses validation
MAE; test data informs no choice.

\paragraph{Feature groups.} \emph{latent}: member-0 belief state $[h_t, z_t]$. \emph{apriori}:
one-step-ahead entropy and ensemble Jensen--Shannon divergence, prior/posterior standard
deviations and their KL at the anchor, one-step tracking error over the observed context, and
distance from the anchor latent to the nearest training latent. \emph{imag}: per-horizon
ensemble disagreement, entropy, prior standard deviation and spread of imagined displacement at
$h \in \{1,2,5,10,20,30\}$ plus summaries, and the horizon at which imagined rollouts first
diverge by one cell (the online-truncation baseline). \emph{actions}: counts and prefix sums of
the planned action sequence. \emph{hand}: object-type counts in the current view, visible free
cells ahead, unseen cells ahead, carrying and door flags, direction. Variants used in the paper
are \texttt{full} (all groups), \texttt{no\_latent}, \texttt{uncertainty\_only}
(\emph{apriori}+\emph{imag}), \texttt{strict\_apriori} (all except \emph{imag}),
\texttt{hand\_actions\_only} and \texttt{actions\_only}. In the cross-architecture experiment
only \texttt{no\_latent} is used, because its column \emph{names} are identical across
architectures while total feature dimensionality is not ($634$, $346$, $362$); the code asserts
name equality before any cross-architecture comparison.

\subsection{Statistics}

Confidence intervals use a bootstrap that resamples whole \emph{episodes}, not anchors, because
anchors within an episode are strongly dependent \citep{efron1993bootstrap}. We use $2000$
resamples and report $95\%$ percentile intervals. Comparisons between methods are paired: the
same resampled episode set is used for both methods and the interval is computed on the
difference. Reported statistics are MAE, RMSE, Spearman $\rho$, Pearson $r$, Kendall $\tau$ and
$R^2$; the main text reports MAE, $\rho$ and $R^2$. Calibration uses the expected calibration
error of $P(\Hstar \ge h)$ in ten equal-width probability bins, discarding bins with fewer than
$20$ states.

The Claim-A ceiling uses isotonic regression \citep{barlow1972isotonic} with the direction
chosen automatically, fitted on the evaluation split. This is deliberately optimistic: it is
the minimum of \eqref{eq:ceiling} attainable by any monotone map, with no generalisation gap.

\subsection{Censoring}

$\Hstar$ is right-censored at $H_{\max} = 30$. On A-test $20.7$--$21.7\%$ of anchors are at the
ceiling across seeds. Consequences: MAE is inflated relative to an uncensored target; the upper
tail of $\Hstar$ is unidentified; and $R^2$ is computed against a censored variance. All
methods including baselines and the ceiling face the identical censoring, so comparisons
\emph{between} methods remain valid, but absolute values should not be read as estimates of an
uncensored horizon.

\subsection{Compute}

The main run trained $15$ world models ($3$ seeds $\times$ $K=5$) in $4.4$ hours on a single
Tesla P100. The cross-architecture study added $9$ models. All follow-up experiments (H5, H6,
H8, H9) reuse stored features or checkpoints and require no world-model retraining. Total:
$24$ world models. Runs were executed on Kaggle with PyTorch $2.5.1{+}$cu121; the preinstalled
build did not support the allocated P100's compute capability, so a Pascal-compatible build is
installed at kernel start.

\section{Additional Results}
\label{app:additional}

\subsection{Full hypothesis ledger}

Table~\ref{tab:ledger} in the main text lists all nine hypotheses. Two procedural notes belong
with it.

First, in two places our scoring code and our written decision rule diverged. In H8 the
pre-registration assigned \textsc{not supported} only when the first prediction failed and
\textsc{partial} when the first two held, leaving the case ``first holds, second fails''
unspecified; the code defaulted to \textsc{not supported} and we report that rather than
choosing a label after seeing the outcome. In H9 the pre-registration states that the verdict
is evaluated only over configurations passing the validity gate, which would make two
predictions uninformative, while the code scored one of them as a failure. Both paths reach the
same verdict, so nothing changes, but we record the discrepancy.

Second, the ensemble-size ablation (seed 0) shows the effect does not require an ensemble:

\begin{center}\small
\begin{tabular}{cccc}
\toprule
$K$ & monitor MAE $\downarrow$ & $\rho$ $\uparrow$ & ceiling MAE \\
\midrule
1 & 6.908 & 0.533 & 8.897 \\
3 & 6.582 & 0.547 & 7.809 \\
5 & 6.504 & 0.559 & 7.583 \\
\bottomrule
\end{tabular}
\end{center}

At $K=1$ there is no ensemble disagreement at all, only single-model entropy and latent
statistics, and the monitor still beats its ceiling --- by the largest margin of the three.
Ensembling improves absolute accuracy and improves the ceiling faster than it improves the
monitor.

\subsection{Size-matched transfer control}

Table~\ref{tab:transfermatched} repeats the transfer matrix with every regime subsampled to the
A-only training-anchor count ($22{,}200$), so all regimes see equal data.

\begin{table}[h]
\centering\small
\input{tables/tab_transfer_matched}
\caption{Size-matched control for Table~\ref{tab:transfer}. Spearman $\rho$, mean over three
seeds; italic entries are configurations seen in training.}
\label{tab:transfermatched}
\end{table}

The critical cell $\mathrm{A{+}B{+}C}\to$D gives MAE $3.622$ and $\rho = 0.762$, against
$3.584$ and $0.766$ unmatched. The effect is therefore not explained by sample count alone. We
do not claim that diversity is its only cause: severity of the donor shift covaries with its
presence, and we did not manipulate the two independently.

\subsection{Rejected instruments}

Two configurations failed the pre-registered validity gate (a new configuration counts as an
instrument only if $\Hstar$ is at least $40\%$ below A's \emph{and} the no-exposure monitor
transfers to it with $\rho < 0.35$):

\begin{center}\small
\begin{tabular}{llcc}
\toprule
& perturbation & $\Hstar$ & no-exposure $\rho$ \\
\midrule
F  & obs.\ noise $0.15$ & 15.23 & $+0.472$ \\
F2 & obs.\ noise $0.70$ & 12.72 & $+0.359$ \\
F3 & action dropout $0.20$ & \phantom{0}3.67 & $+0.238$ \\
\bottomrule
\end{tabular}
\end{center}

Observation noise never bites: even at $70\%$ cell corruption $\Hstar$ falls only from $15.5$
to $12.7$. An action-conditioned model with sixteen steps of belief context does not depend
heavily on the corrupted channel. F and F2 are reported as invalid instruments rather than as
evidence about transfer; predictions targeting them are uninformative by rule. The severity
ladder used to select F2 and F3 is in \texttt{results/aleatoric/calibration.json}.

A useful control follows from this. The regime $\mathrm{A{+}F2}$ --- exposure to a
configuration that barely degrades the model --- teaches nothing transferable: $\rho = 0.226$
on B and $0.087$ on D, at or below the no-exposure baseline. Exposure helps only when the donor
genuinely degrades the world model.

\subsection{Planning experiment (weak test)}

The adaptive-planning arm sets the imagination horizon of a random-shooting model-predictive
controller to the monitor's prediction, clipped to $[1, 30]$, and is charged for the probe
rollout used to produce that prediction ($K \cdot H_{\max}$ imagined states per decision).
Candidate action sequences are drawn from the same forward-biased distribution used to collect
training data. Targets require a detour: shortest-path distance must exceed Manhattan distance.
Sixty layouts, paired across arms.

\begin{center}\small
\begin{tabular}{lccc}
\toprule
arm & success & imagined states / decision & mean $H$ \\
\midrule
fixed $H{=}3$  & $0.917 \pm 0.036$ & \phantom{0}288 & \phantom{0}3.0 \\
fixed $H{=}5$  & $0.933 \pm 0.032$ & \phantom{0}480 & \phantom{0}5.0 \\
fixed $H{=}10$ & $0.767 \pm 0.055$ & \phantom{0}960 & 10.0 \\
fixed $H{=}20$ & $0.517 \pm 0.065$ & 1920 & 20.0 \\
fixed $H{=}30$ & $0.367 \pm 0.062$ & 2880 & 30.0 \\
adaptive       & $0.517 \pm 0.065$ & 2637 & 25.9 \\
\bottomrule
\end{tabular}
\end{center}

Success falls monotonically with horizon, so the compute--performance frontier has no interior
optimum and an adaptive rule has nothing to find. The detour manipulation did not take: mean
detour was $0.13$ cells, with $57$ of $60$ layouts falling through to the least restrictive
tier. We therefore report this as a weak test of downstream utility rather than a refutation.
The adaptive arm's specific failure --- selecting $H \approx 26$ where $3$--$5$ is optimal ---
is the anti-conservative bias of \S\ref{sec:ood} reappearing under the state distribution that
model-predictive control itself induces.

\subsection{Information-boundary verification}

The monitor may use no observation, state, or error at time $>t$. This is verified
executably rather than asserted: every field of the evaluation window strictly after the anchor
is replaced with random values, the features are recomputed, and all feature columns must be
bit-identical while the target must change. Corrupting the future leaves the features
unchanged and moves mean $\Hstar$ from $6.45$ to $0.01$. The planned action sequence is
deliberately preserved by the corruption, since a planner genuinely knows it.


\section{Deviations from Pre-Registration}
\label{app:deviations}

Ten deviations are logged in \texttt{DEVIATIONS.md} with their timing relative to test
evaluation. Six occurred before any test data was examined:

\begin{itemize}\itemsep1pt
\item[D1] Configurations changed after a geometry survey that used no neural network: the
original set left the agent travelling $\approx 2.7$ cells in $30$ steps while seeing $4$ cells
ahead, so the planned path never left the visible region and $\Hstar$ was degenerate.
\item[D2] Collection policy changed from a BFS explorer to a forward-biased one, for the reason
given in \S\ref{app:details}.
\item[D3] Splits restructured to give the monitor its own world-model-held-out split.
\item[D4] Pilot test seeds moved outside the main test range after an early diagnostic reported
aggregate accuracies on an overlapping range; that diagnostic informed no threshold and its
models and data were discarded.
\item[D5] Two analysis additions: the privileged-state comparison, and selection of the monitor
variant on validation MAE.
\item[D6] Planning task made non-myopic, and the adaptive arm charged for its probe rollout.
\end{itemize}

Four are post-main-run and scoped to follow-ups: D7 records that all reported numbers come from
Kaggle and that a briefly started local run was terminated without producing any result; D8
records the planning re-run and its residual limitation; D9 records that the H5 protocol was
frozen blind; D10 records the size-matched control and the new $60/20/20$ partitions, and notes
that consequently H6 numbers are not directly comparable to full-split numbers elsewhere.

\bibliographystyle{plainnat}
\bibliography{../references}

\end{document}
